% ============================================================================
%  abstract.tex — Résumé / Abstract Template
%  Replace the placeholder text below with your project's summary.
% ============================================================================

% --- French Abstract ---
\chapter*{Résumé}
\addcontentsline{toc}{chapter}{Résumé}

\vspace{0.5cm}

Dans le cadre du module de \textbf{Gestion de Projets Informatiques} dispensé à l'École Nationale des Sciences Appliquées d'Al Hoceima, ce projet vise à concevoir, implémenter et valider une \textbf{architecture Kappa} appliquée à une plateforme e-commerce fictive.

\vspace{0.3cm}

Face à la complexité croissante des architectures de données traditionnelles --- notamment l'architecture \textbf{Lambda} qui impose la maintenance de deux couches distinctes de traitement (batch et temps réel) ---, l'architecture \textbf{Kappa}, proposée par Jay Kreps en 2014, offre une alternative élégante reposant sur un \textbf{unique flux de traitement continu}. Ce projet démontre concrètement cette approche en déployant une infrastructure complète composée d'\textbf{Apache Kafka} (log d'événements), \textbf{Apache Flink} (référence de traitement de flux) et \textbf{Docker} (conteneurisation).

\vspace{0.3cm}

Le système mis en place traite en temps réel des événements comportementaux simulés (vues, ajouts au panier, achats, abandons) à travers trois topologies spécialisées : le \textbf{filtrage des bots}, l'\textbf{agrégation statistique par fenêtres temporelles} et la \textbf{détection d'abandon de panier} avec alertes marketing. La validation du système repose sur le mécanisme de \textbf{replay} --- pierre angulaire de l'architecture Kappa --- qui permet de rejouer l'intégralité du journal d'événements pour recalculer les résultats et confirmer leur parfaite cohérence avec le traitement temps réel.

\vspace{0.3cm}

Le projet a été conduit selon une méthodologie \textbf{Agile (Scrum + Kanban)}, structuré en 4 sprints, avec une gestion rigoureuse des risques et une planification par \textbf{Planning Poker}. L'ensemble du processus a été suivi via \textbf{Taiga} (Kanban) et \textbf{GanttProject} (planification temporelle).

\vspace{0.5cm}

\noindent\textbf{Mots-clés :} Architecture Kappa, Apache Kafka, Traitement de flux, Temps réel, E-commerce, Rétention, Replay, Docker, Gestion de projet Agile.

\vfill
